% Part: normal-modal-logic
% Chapter: axioms-systems
% Section: logics-proofs

\documentclass[../../../include/open-logic-section]{subfiles}

\begin{document}

\olfileid{nml}{prf}{prf}

\olsection{推导与模态系统}

我们首先定义正规模态逻辑中的推导。粗略而言，推导是一个公式序列，其中每个元素要么是若干\emph{公理}之一的代入实例，要么通过某几条推理规则之一由前面的元素得出。对于正规模态逻辑，所有重言式实例\iftag{prvDiamond}{、\Ax{K} 和~\Dual{}}{以及~\Ax{K}}均算作公理。这就得到了模态系统~$\Log{K}$，即最小的正规模态逻辑。不过，我们有时希望增加额外公理以得到其他系统。推理规则则始终是肯定前件~\MP{} 和必然化规则~\Nec。

\begin{defn}
  给定一个模态系统 $\Log{K} !A_1 \dots !A_n$ 和一个公式~$!B$，我们称 $!B$ 在 $\Log{K} !A_1 \dots
  !A_n$ 中是\emph{可推导的}，记作 $\Log{K} !A_1 \dots !A_n \Proves !B$，当且仅当存在公式 $!C_1$, \dots, $!C_k$ 使得 $!C_k = !B$，且每个 $!C_i$ 要么是一个重言式实例，要么是 $\Ax{K}$、\iftag{prvDiamond}{ $\Dual$、}{} $!A_1$, \dots, $!A_n$ 中之一的实例，要么是通过规则 \MP{} 或 \Nec 由前面的公式得出。
\end{defn}

下述命题使得我们能够通过展示~$!B$ 的一个 $\Sigma$-推导来证明 $!B \in \Sigma$。

\begin{prop}
  $\Log{K} !A_1 \dots !A_n = \Setabs{!B}{\Log{K} !A_1 \dots !A_n \Proves !B}$.
\end{prop}

\begin{proof}
  我们对推导的长度进行归纳，以证明 $\Setabs{!B}{\Log{K} !A_1 \dots !A_n \Proves !B} \subseteq
  \Log{K} !A_1 \dots !A_n$。

  若~$!B$ 的推导长度为~$1$，则它仅包含一个公式。该公式不可能通过 \MP{} 或 \Nec 由前面的公式得出，因此必定是一个重言式实例，或是 \Ax{K}\iftag{prvDiamond}{、$\Dual$}{} 或 $!A_1$, \dots,~$!A_n$ 中之一的实例。但 $\Log{K}!A_1\dots!A_n$ 也包含这些公式，故有 $!B \in \Log{K}!A_1 \dots !A_n$。

  若 $!B$ 的推导长度~$> 1$，则 $!B$ 还可能是由推导中非最后一行的公式通过 \MP{} 或 \Nec{} 得到的。若 $!B$ 由 $!C$ 和 $!C \lif !B$ 通过 \MP 得出，则由归纳假设，$!C$ 与 $!C \lif !B \in
  \Log{K}!A_1 \dots !A_n$。但每个模态逻辑都对肯定前件封闭，因此 $!B \in \Log{K}!A_1 \dots
  !A_n$。若 $!B \equiv \Box !C$ 由 $!C$ 通过 \Nec 得出，则由归纳假设 $!C
  \in \Log{K}!A_1 \dots !A_n$。但每个正规模态逻辑都对 $\Nec$ 封闭，因此 $!B \in
  \Log{K}!A_1\dots!A_n$。

  相反的包含关系可通过证明 $\Sigma = \Setabs{!B}{\Log{K} !A_1 \dots !A_n \Proves !B }$ 是一个包含 $!A_1$, \dots, $!A_n$ 的所有实例的正规模态逻辑，并根据定义 $\Log{K} !A_1 \dots !A_n$ 是这类逻辑中最小的一个而得证。
  \begin{enumerate}
    \item 每个重言式~$!B$ 都是重言式实例，因此 $\Log{K}!A_1\dots!A_n \Proves !B$，故 $\Sigma$ 包含所有重言式。
    \item 若 $\Log{K}!A_1\dots!A_n \Proves !C$ 且 $\Log{K}!A_1\dots!A_n \Proves !C \lif !B$，则 $\Log{K}!A_1\dots!A_n \Proves !B$：将 $!C$ 的推导与 $!C \lif !B$ 的推导合并，并添加一行~$!B$。最后一行的依据是 \MP{}。所以 $\Sigma$ 对肯定前件封闭。
    \item 若 $!B$ 有一个推导，则~$!B$ 的每个代入实例也有一个推导：对该推导中的每个公式施行代入。（练习：通过对推导长度归纳，证明结果也是一个正确的推导。）因此 $\Sigma$ 对一致代入封闭。（至此我们已确立 $\Sigma$ 满足模态逻辑的所有条件。）
    \item 我们有 $\Log{K}!A_1\dots!A_n \Proves \Ax{K}$，故 $K \in \Sigma$。
    \tagitem{prvDiamond}{我们有 $\Log{K}!A_1\dots!A_n \Proves \Dual$，故 $\Dual \in \Sigma$。}{}
    \item 若 $\Log{K}!A_1\dots!A_n \Proves !C$，则附加行~$\Box !C$ 的依据是~\Nec。因此，$\Sigma$ 对~\Nec 封闭。故 $\Sigma$ 是正规模态逻辑。
  \end{enumerate}
\end{proof}

\end{document}